%======================================================
% Conference Booklet / Book of Abstracts — adapted from AMCOS_booklet
% https://github.com/maximelucas/AMCOS_booklet (GPL-3.0)
% Original authors: Maxime Lucas, Pau Clusella
%
% Simplified from upstream: dropped the online/print dual-compile
% machinery (\ifOnline switch, two separate abstract sets under
% abstracts/tex/) and the cover page, which upstream includes as an
% external PDF ("our cover was produced with canva.com" — not
% redistributable without that design tool). This starter draws its own
% simple title page in LaTeX instead, and ships ONE version with real
% (invented) abstract content in place of upstream's Jane-Austen-style
% filler text.
%======================================================
\documentclass[openany, parskip=full, 12pt, a4paper]{scrbook}

\usepackage[utf8]{inputenc}
\usepackage[T1]{fontenc}
\usepackage{xspace}
\usepackage{microtype}
\usepackage{cmbright}
\usepackage{color}
\usepackage[table]{xcolor}
\usepackage{scrhack}
\usepackage{tikz}
\usepackage{graphicx}
\usepackage{float}
\usepackage{multicol}
\usepackage{longtable}
\usepackage{ifthen}
\usepackage[english]{babel}
\usepackage[hidelinks]{hyperref}

\pagestyle{plain}
\setcounter{secnumdepth}{-2}

%======================================================
% Adapted from AMCOS_booklet (GPL-3.0)
% https://github.com/maximelucas/AMCOS_booklet
% Original authors: Maxime Lucas, Pau Clusella
%=======================================================

\definecolor{myorange}{RGB}{255,117,40}
\definecolor{mygray}{RGB}{164, 168, 172}
\definecolor{mywhite}{RGB}{235, 238, 231}
\definecolor{myblue}{RGB}{52, 115, 116}

\newcommand{\primarycolor}{myblue}
\newcommand{\secondarycolor}{mywhite}

%----------------------------------------
% talk abstract environment
%----------------------------------------
\newenvironment{abstract}[4] %{title}{author}{affiliation}{tag}
{\filbreak

	{\large \bfseries #1}

	{\bfseries \itshape #2,} \textcolor{mygray}{#3} \hfill {#4}

}
{}

%----------------------------------------
% poster abstract command
%----------------------------------------
\newcommand{\poster}[3] %{title}{author}{affiliation}
{\filbreak

	{\bfseries \large #1} \\
	\tab #2, \textit{#3}

}

\newcommand\tab[1][1cm]{\hspace*{#1}}

%----------------------------------------
% tags for talk type (coloured circle in abstracts)
%----------------------------------------

\newcommand{\KLcolor}{myorange!45}
\newcommand{\IScolor}{myblue!65}
\newcommand{\CTcolor}{white}

\newcommand{\KLtag}{\tikz[baseline={([yshift=-.8ex]current bounding box.center)}]  \node[circle, inner sep=2pt, minimum size=0.5em, color=black, fill=\KLcolor]{\small \bfseries KL};}
\newcommand{\IStag}{\tikz[baseline={([yshift=-.8ex]current bounding box.center)}]  \node[circle, inner sep=2pt, minimum size=0.5em, color=black, fill=\IScolor]{\small \bfseries IS};}
\newcommand{\CTtag}{\tikz[baseline={([yshift=-.8ex]current bounding box.center)}]  \node[circle, inner sep=2pt, minimum size=0.5em, color=black, fill=\CTcolor]{\small \bfseries CT};}

%
% PAGE LAYOUT: vertical colour bar on the outer edge of each page
%
\usepackage{etoolbox}
\usepackage[scale=1,angle=0,opacity=1]{background}
\backgroundsetup{contents={}}

\AddEverypageHook{%
\ifthenelse{\isodd{\thepage} \AND \thepage>1}{%
	\backgroundsetup{color=\secondarycolor,position=current page.south east,
		nodeanchor=south east,contents={\rule{10pt}{0.66\paperheight}}}
	}{}
\ifthenelse{\NOT \isodd{\thepage}}{%
	\backgroundsetup{color=\secondarycolor,position=current page.south west,
		nodeanchor=south west,contents={\rule{10pt}{0.66\paperheight}}}
	}{}
\BgMaterial}

%---------------------------------------------------
% chapter heading style — coloured bar beside the title
%---------------------------------------------------
\newdimen\mybarpadding
\mybarpadding=1.5em\relax

\RedeclareSectionCommand[
    afterskip=4em plus 1pt minus 1pt,
    beforeskip=-1pt,
    level=0,
    toclevel=0
]{chapter}

\setkomafont{chapter}{\normalfont\normalsize\bfseries\Huge}

\makeatletter
\renewcommand{\chapterlinesformat}[3]{%
  \ifthispageodd{%
    \hfill%
    \raisebox{-0.2em}{\makebox[0pt][r]{\textcolor{\primarycolor}{\rule{\paperwidth}{1em}}}}%
    \hspace{\mybarpadding}%
	\mbox{#3}%
  }{%
      \mbox{#3}%
      \hspace{\mybarpadding}%
      \raisebox{-0.2em}{\makebox[0pt][l]{\textcolor{\primarycolor}{\rule{\paperwidth}{1em}}}}%
  }%
}
\makeatother

%
% TIMETABLE COLOURS AND STYLES
%
\newcommand{\tbg}{gray}
\newcommand{\tfg}{white}
\newcommand{\tbc}{gray!25}

\newcommand{\tablebreak}[2]{\rowcolor{\tbc} #1 &  \multicolumn{4}{c|}{\bfseries #2} \\ \hline}
\newcommand{\eventtype}[2]{#1& \multicolumn{4}{c|}{\cellcolor{\tbg}\color{\tfg}\bfseries #2} \\ \hline}

\newcolumntype{L}[1]{>{\raggedright\let\newline\\\arraybackslash\hspace{0pt}}m{#1}}
\newcolumntype{C}[1]{>{\centering\let\newline\\\arraybackslash\hspace{0pt}}m{#1}}
\newcolumntype{R}[1]{>{\raggedleft\let\newline\\\arraybackslash\hspace{0pt}}m{#1}}

\newcommand{\IS}[5]{#1 &\cellcolor{\IScolor}IS&{\bfseries#2}\newline #4&&#5 \\ \hline}
\newcommand{\CT}[5]{#1 &\cellcolor{\CTcolor}CT&{\bfseries#2}\newline #4&&#5 \\ \hline}
\newcommand{\KL}[5]{#1 &\cellcolor{\KLcolor}KL&{\bfseries#2}\newline #4&&#5 \\ \hline}


\begin{document}

% TITLE PAGE
%--------------------------------------------------------------------
\begin{titlepage}
\centering
\vspace*{3cm}
{\Huge \bfseries \textcolor{myblue}{Complex Networks in Motion}}\\[1em]
{\Large 4th International Workshop on Dynamical Networks}\\[2em]
{\large Book of Abstracts}\\[3em]
{\large 14--16 September 2026 \\ Coastal University, Porthaven}
\vfill
{\small The open-source \LaTeX{} template used to typeset this booklet,
AMCOS\_booklet, is available at
\url{https://github.com/maximelucas/AMCOS_booklet}.}
\end{titlepage}

\tableofcontents

% ABOUT
%---------------------------------------------------------------------
\chapter{About}

\section{The workshop}
Complex Networks in Motion brings together researchers studying how
networked systems evolve over time — from neural populations to power
grids — with a focus on methods that combine dynamical-systems theory
with network inference from data.

\section{Organizing committee}
\begin{center}
\begin{tabular}{lll}
Amara Okafor & Diego Fontan & Ines Larsen \\
Bo Lindqvist & Farah Haddad & Jian Wu \\
Camille Dubois & Grace Muthoni & Karim El-Sayed
\end{tabular}
\end{center}

% TIMETABLE
%---------------------------------------------------------------------
\chapter{Timetable}

KL: Keynote Lecture, IS: Invited Speaker, CT: Contributed Talk.

\begin{center}
\renewcommand{\arraystretch}{1.3}
\begin{longtable}{|C{0.15\linewidth}| C{0.04\linewidth}|  C{0.3\linewidth} C{0.0\linewidth} C{0.4\linewidth}|}
\hline
\eventtype{9:00}{Registration \& Coffee}
\KL{9:30}{Sofia Ramirez}{}{Coastal University}{Synchronisation in adaptive networks}
\IS{10:30}{Tomasz Nowak}{}{Baltic Institute}{Inferring coupling from time series}
\tablebreak{11:15}{Coffee Break}
\CT{11:30}{Priya Anand}{}{Delta Tech}{A spectral test for hidden community structure}
\CT{12:00}{Liu Wei}{}{Northshore Institute}{Delay-coupled oscillators on random graphs}
\end{longtable}
\end{center}

% TALKS
%---------------------------------------------------------------------
\chapter{List of Abstracts --- Talks}

\section{Monday 14th}

\begin{abstract}{Synchronisation in adaptive networks}{%
	S. Ramirez}{%
	Coastal University, Porthaven}{%
	\KLtag}
We study networks whose coupling weights adapt in response to the
synchrony of the nodes they connect, a mechanism observed in both neural
plasticity and power-grid load balancing. We derive a reduced set of
equations for the joint evolution of phase and weight dynamics, and show
that adaptive coupling can stabilise partial synchrony states that are
unstable under fixed coupling. Numerical continuation confirms the
prediction on networks of up to 200 Kuramoto oscillators.
\end{abstract}

\begin{abstract}{Inferring coupling from time series}{%
	T. Nowak}{%
	Baltic Institute, Gdynia}{%
	\IStag}
Reconstructing the interaction structure of a dynamical network from
observed trajectories alone is a central inverse problem in network
science. We compare three inference methods — transfer entropy,
Granger causality, and a phase-response-curve-based estimator — on
synthetic data with known ground truth, and report substantial
sensitivity to observation noise and sampling rate that is rarely
discussed in application papers.
\end{abstract}

\section{Monday 14th, continued}

\begin{abstract}{A spectral test for hidden community structure}{%
	P. Anand}{%
	Delta Tech, Bengaluru}{%
	\CTtag}
We propose a spectral statistic, computed from the normalized graph
Laplacian, that distinguishes networks with planted community structure
from Erdős--Rényi null models at sample sizes where modularity-based
tests lose power. The statistic requires no prior estimate of the
number of communities.
\end{abstract}

\begin{abstract}{Delay-coupled oscillators on random graphs}{%
	L. Wei}{%
	Northshore Institute, Shanghai}{%
	\CTtag}
Transmission delays between coupled oscillators are unavoidable in
physically extended networks. We characterise how delay interacts with
random graph topology to produce multistability between in-phase and
splay states, and give an analytic condition, in terms of the graph's
spectral gap, for the in-phase state to remain the unique attractor.
\end{abstract}

% POSTERS
%------------------------------------------------------------------
\chapter{List of Posters}

\vspace{-1em}

\poster{Epidemic spreading with heterogeneous recovery rates}{J. Almeida}{Riverside College}
\poster{Community detection under partial observability}{H. Park}{Northshore Institute}
\poster{A dataset of intercity mobility networks, 2018--2025}{M. Haile}{Delta Tech}

% USEFUL INFO
%------------------------------------------------------------------
\chapter{Useful Information}

\textbf{Venue.} All sessions take place in the Bayview Lecture Hall,
Coastal University.

\textbf{Wi-Fi.} Network \texttt{CU-Guest}, no password required.

\textbf{Lunch.} Served daily in the adjacent Refectory, 12:30--13:30.

\end{document}
